\section{课程间知识点的（相互）关联}
\subsection*{话题 2-A：极限与积分交换次序}
这里将一致收敛性、Riemann 积分、反常积分及 Lebesgue 控制收敛定理放在同一框架下比较。

\thm{7.2.1}若 $f_n\in C[a,b]$ 且 $f_n\rightrightarrows f$，则 $f$ 连续，且
\[\lim_{n\to\infty}\int_a^b f_n(x)\dd x=\int_a^b f(x)\dd x.\]
\thm{7.2.2}若 $f_n$ 在 $[a,b]$ 上 Riemann 可积且一致收敛于 $f$，则 $f$ 也 Riemann 可积，上述等式成立。
\pf 由上、下积分和 $\|f_n-f\|_\infty\to0$，分别给出 $\overline{\int}f\le\lim\int f_n$ 与 $\underline{\int}f\ge\lim\int f_n$；二者相等。
\thm{7.2.3（控制收敛）}设 $E\subset\R^n$ Lebesgue 可测，$f_k,g$ 可积，$|f_k|\le g$ 几乎处处，且 $f_k\to f$ 几乎处处。则 $f$ 可积且
\[\lim_{k\to\infty}\int_E f_k(x)\dd x=\int_E f(x)\dd x.\]
\par 一致收敛是很强的假设。Lebesgue 控制收敛定理则通过可积的控制函数，允许几乎处处收敛；它同时覆盖有界区域和无界区域。Riemann 积分的极限函数即使有界也可能不再 Riemann 可积，因此不能把这个结论直接写成不加条件的 Riemann 版。
\thm{7.2.4（Arzelà 有界收敛）}设 $E\subset\R^n$ 有界且可求体积。若 $f_k$ 在 $E$ 上一致有界、Riemann 可积，$f_k(x)\to f(x)$，且 $f$ 亦 Riemann 可积，则 $\int_E f_k\to\int_E f$。
\par 定理 7.2.4 与定理 7.2.3 的关系值得注意：若把 Riemann 可积函数视作 Lebesgue 可积函数，则一致有界可取常值函数作控制。困难只在于极限的 Riemann 可积性需要另加假设；Lebesgue 理论没有这个障碍。

\thm{7.2.5}若连续函数列 $f_n$ 在 $[a,\infty)$ 上内闭一致收敛于 $f$，且 $\int_a^\infty f_n$ 关于 $n$ 一致收敛，则 $\int_a^\infty f$ 收敛并可交换极限与反常积分。
\thm{7.2.6}例 7.2.5 中可将“连续”减弱为：对每个 $A>a$，$f_n$ 在 $[a,A]$ 上 Riemann 可积；其余条件不变，结论仍成立。
\thm{7.2.7}若 $f_n$ 在 $[a,\infty)$ 收敛于 $f$，每个 $[a,A]$ 上的 $f_n$ 一致有界且 Riemann 可积，$f$ 亦可积，且 $\int_a^\infty f_n$ 关于 $n$ 一致收敛，则仍可交换极限和反常积分。
\pf 证明这三条结论的共同结构是“先截断，再控制尾部”。固定 $A>a$，在 $[a,A]$ 上分别应用定理 7.2.2 或 7.2.4；然后把
\[\int_a^\infty f_n-\int_a^\infty f
=\int_a^A(f_n-f)+\int_A^\infty f_n-\int_A^\infty f\]
分为局部误差和两个尾积分。局部误差通过 $n\to\infty$ 消失；反常积分的一致收敛使尾积分对所有充分大的 $n$ 同时小于给定的 $\varepsilon$。这种顺序不可颠倒：只对每个固定 $n$ 分别有尾积分收敛，仍可能在 $n\to\infty$ 时出现“质量逃向无穷”。
\thm{7.2.8}若 $f_n$ 为 $[a,\infty)$ 上的 Lebesgue 可测函数，$|f_n|\le g$ 几乎处处，其中 $g$ 可积，且 $f_n\to f$ 几乎处处，则
\[\lim_{n\to\infty}\int_a^\infty f_n(x)\dd x=\int_a^\infty f(x)\dd x.\]
\par 与前面的 Riemann 型结论相比，这里同一控制函数 $g$ 同时处理局部收敛与无穷远处的尾部。形式上只需 $|f_n|\le g$，无需另外单列“反常积分一致收敛”的条件。
例如，$0\le\int_0^1 x^n/(1+x)\dd x\le1/(n+1)$，而
\[\lim_{n\to\infty}\int_0^n\frac{x}{(1+x/n)^n}\dd x
=\int_0^\infty xe^{-x}\dd x=\Gamma(2)=1.\]
后一式可将被积函数延拓为 $x(1+x/n)^{-n}\mathbf1_{[0,n]}(x)$，再用可积控制函数及控制收敛定理。
\par 对后一个极限也可以直接证明控制函数的存在：当 $n\ge4$ 且 $0\le x\le n$ 时，二项式展开给出
\[\left(1+\frac{x}{n}\right)^n\ge1+\binom n3\frac{x^3}{n^3}\ge1+\frac{x^3}{24}.\]
因此 $x(1+x/n)^{-n}\mathbf1_{[0,n]}(x)\le24x/(24+x^3)$；后者在 $[0,\infty)$ 可积。与此相比，第一式只需要 $x^n/(1+x)\le x^n$，无需使用抽象定理。

\subsection*{话题 2-E：代数与分析的相互影响}
\ex{7.2.32}证明：线性变换特征值的几何重数不超过其代数重数。
\pf 若 $\dim\ker(A-\lambda I)=r$，先取 $r$ 个线性无关的特征向量扩充为一组基。该基下矩阵呈分块上三角形，左上角为 $\lambda I_r$，故特征多项式含因子 $(t-\lambda)^r$。
\par 原书进一步用矩阵 $\begin{psmallmatrix}2&1\\0&2\end{psmallmatrix}$ 及扰动 $\begin{psmallmatrix}2+\varepsilon&1\\0&2\end{psmallmatrix}$ 说明：当 $\varepsilon\to0$ 时，两条特征方向合拢，解释了 Jordan 块对应的几何重数下降。
\par 对 $A_\varepsilon=\begin{psmallmatrix}2+\varepsilon&1\\0&2\end{psmallmatrix}$，两个特征值为 $2+\varepsilon$ 与 $2$。可分别取特征向量 $(1,0)^\mathsf T$ 与 $(1,-\varepsilon)^\mathsf T$。当 $\varepsilon\to0$，这两个方向趋于同一方向；从几何上看，原来的二维特征空间退化成一条直线。对三阶 Jordan 块可作同样的扰动解释：一条长度为三的链被“拆开”成三个相近而不同的特征方向。
\ex{7.2.33}类似地，讨论三个三阶上三角矩阵：一条长度为三的 Jordan 链、长度为二的链加一条独立特征向量，以及 $2I_3$。它们的代数重数均为三，几何重数分别为 $1,2,3$。
\ex{7.2.34}求 $y''-2y'+y=0$ 的通解。
\sol 特征方程 $(\lambda-1)^2=0$。一个解为 $e^x$。对相邻不同根的解 $e^{(1+\delta_1)x},e^{(1+\delta_2)x}$ 作差商，根合并时趋于 $xe^x$，故
\[y(x)=C_1e^x+C_2xe^x.\]
\par 这里的差商说明重根对应的解 $xe^x$ 从何而来。若把方程稍微扰动，使特征方程出现两个不同根 $1+\delta_1$ 和 $1+\delta_2$，则
\[\frac{e^{(1+\delta_1)x}-e^{(1+\delta_2)x}}{\delta_1-\delta_2}\longrightarrow xe^x.\]
因此额外解并非凭空猜测，而是相邻特征根的极限。

\subsection*{话题 2-G：微分方程与分析估计}
\ex{7.2.42}若 $\lim_{x\to\infty}(f'(x)+2f(x))=\ell$，则 $\lim_{x\to\infty}f(x)=\ell/2$。
\pf 因 $f=e^{-2x}(e^{2x}f)$，分子、分母同趋无穷时应用 L'Hospital 法则，极限为 $(f'+2f)/2\to\ell/2$。
\ex{7.2.43}设 $f,g\in C(\R)$，$g>0$，$\int_0^\infty g=+\infty$，$f(x)/g(x)\to0$，且 $y'+gy=f$。证明 $y(x)\to0$。
\pf 令 $G(x)=\int_0^xg(t)\dd t$，则 $G\to\infty$。由 $y=e^{-G}(e^Gy)$，应用 L'Hospital 法则：
\[\lim_{x\to\infty}y(x)=\lim_{x\to\infty}\frac{(e^Gy)'}{(e^G)'}=\lim_{x\to\infty}\frac{f(x)}{g(x)}=0.\]
\ex{7.2.44}设 $\alpha>1$，证明不存在 $[0,\infty)$ 上的正可导函数 $f$ 满足 $f'(x)\ge f(x)^\alpha$。
\pf 若存在，则 $\bigl(x+f^{1-\alpha}(x)/(\alpha-1)\bigr)'\le0$，故 $x\le f(0)^{1-\alpha}/(\alpha-1)$，与任意大的 $x$ 矛盾。
\ex{7.2.45}设 $m,n$ 为正整数、$k\ne0$。证明不存在定义在整个 $\R$ 上且满足
\[y'(x)=k y(x)^{2n}+x^{2m-1}\]
的函数。
\pf 令 $\varepsilon=\operatorname{sgn}k$、$f(x)=y(\varepsilon x)$，则 $f'(x)=|k|f(x)^{2n}+x^{2m-1}$。在充分大的正半轴上 $f$ 严格增长并为正，最终满足 $f'\ge c f^{2n}$；由例 7.2.44 的爆破估计得矛盾。
\ex{7.2.46}证明不存在 $\R$ 上连续可微的 $f$ 满足 $f'=f\circ f$ 且 $f(0)>0$。
\pf 先证 $f$ 在 $[0,\infty)$ 上为正。若有首个正零点 $\xi$，则 $f>0$ 于 $[0,\xi)$，而 $f'(\xi)=f(f(\xi))=f(0)>0$，与首次从正值降到零矛盾。于是 $f(f(x))>0$，故 $f$ 在正半轴严格递增，且
\[f'(x)=f(f(x))\ge f(f(0))\ge f(0)=:a>0.\]
因此 $f(x)\ge ax$。将其用于 $f(f(x))$，得到 $f'(x)\ge a f(x)$，由积分因子可得 $f(x)\ge f(0)e^{ax}$。再把这个下界用于 $f(f(x))$：
\[\bigl(-e^{-af(x)}\bigr)'=ae^{-af(x)}f'(x)=ae^{-af(x)}f(f(x))\ge a f(0)>0.\]
左侧的原函数始终小于零，却必须至少线性增长，矛盾。
\ex{7.2.47}设 $a>1$，$f:(0,\infty)\to(0,\infty)$ 可微。证明存在 $x_n\to\infty$，使 $f'(x_n)<f(ax_n)$。
\pf 反设在某个尾部 $f'(x)\ge f(ax)>0$。则 $f$ 在该尾部单增，且中值定理给出 $f(ax)-f(x)\ge(a-1)x f(ax)$，取 $x>1/(a-1)$ 矛盾。
\ex{7.2.48}若 $f\in C^1(\R)$，$|f(0)|<1$，$f'=f^2-1$，则 $|f(x)|<1$ 对所有 $x\in\R$ 成立。
\pf 解在 $\pm1$ 处有平衡点。若首次碰到边界，由方程的局部唯一性，与常值解重合，违背 $|f(0)|<1$。
\ex{7.2.49}若 $f\in C^1(\R)$，$|f(0)|<1$，$f'=f^2-2+\sin x$，则 $|f(x)|\le\sqrt3$。
\pf 若在正半轴上首次向上越过 $1$，在首次达到介于初值和 $1$ 的水平时应有非负导数，但 $f^2-2+\sin x<0$，矛盾。类似论证给出正半轴下界 $-\sqrt3$；对负半轴反向使用导数符号得到 $-1<f\le\sqrt3$。
